\chapter{AI And Collaboration Strategy}

\section{AI Feature Set}

The system implements three AI features: rewrite, summarize, and translate. These satisfy the brief's
requirement for at least two AI capabilities. The backend exposes streamed responses so the frontend can
render output progressively. Users can cancel an in-progress generation and can either apply, partially
apply, edit, or reject the final suggestion.

Prompt construction is separated into backend prompt modules and provider abstractions rather than being
hardcoded directly into frontend components. This keeps provider-specific logic in one place and makes it
possible to swap between a deterministic stub provider and a local LM Studio/OpenAI-compatible provider.

\section{AI Under Concurrent Editing}

The main AI/collaboration risk is stale application of suggestions after the source text has already been
changed by another collaborator. The final implementation addresses this explicitly.

When AI is invoked, the system records the source selection snapshot and the exact selected text. When the
user later tries to apply the result, the frontend verifies that the current text inside the target range
still matches that original source. If it does not, the apply is blocked and the user is asked to rerun AI
against the latest text. This prevents silent overwrites and keeps collaboration behaviour consistent under
concurrent editing.

\section{History, Policy, And Quota}

Every AI interaction is logged with document association, action type, prompt label, output preview, and
status. The UI exposes this history per document. AI policy and quota usage are also visible: owners can
enable/disable AI for a document, control which roles may invoke AI, and see quota-related usage metrics.
